% ---------------------------------------------------------------------------
% Trustworthy Agentic Automation for Intent-Based Self-Healing in SDN
% IEEE conference format. Built with tectonic.
%
% Every tabular number comes from final_submission/paper/tables.tex, which is
% generated from final_submission/data/ by
%   final_submission/scripts/make_paper_tables.py
% and every number quoted in the prose is a macro defined in the same file.
% ---------------------------------------------------------------------------
\documentclass[conference]{IEEEtran}

% Tectonic compiles with XeTeX, so the font is selected with fontspec.
% TeX Gyre Termes is the Times-like OpenType face the IEEE template expects.
\usepackage{fontspec}
% TeX Gyre Termes is the Times-like face the IEEE template expects and is
% present on Overleaf; Nimbus Roman is its URW equivalent and is what the
% Tectonic bundle carries.
\IfFontExistsTF{TeX Gyre Termes}{\setmainfont{TeX Gyre Termes}}{\setmainfont{Nimbus Roman}}
\usepackage{amsmath}
\usepackage{amssymb}
\usepackage{algorithm}
\usepackage{algpseudocode}
\usepackage{graphicx}
\usepackage{booktabs}
\usepackage{multirow}
\usepackage{array}
\usepackage[font=footnotesize,labelfont=bf]{xcolor}
\usepackage[hidelinks,breaklinks=true]{hyperref}

% IEEEtran already defines \usepackage{algorithm} support; keep the counters
% predictable.
\newcommand{\ie}{\emph{i.e.}}
\newcommand{\eg}{\emph{e.g.}}
\newcommand{\Repo}{\texttt{final\_submission/}}

% Numbers quoted in the prose, computed from the same data.
\newcommand{\BaselineMean}{17.8}
\newcommand{\BaselineStd}{1.2}
\newcommand{\BaselinePeak}{20.2}
\newcommand{\CongestedMean}{98.1}
\newcommand{\CongestedPeak}{107.3}
\newcommand{\RecoveryAfter}{26.7}
\newcommand{\RecoveryBefore}{97.4}
\newcommand{\PartialAfter}{25.3}
\newcommand{\PartialBefore}{97.4}
\newcommand{\FailedBefore}{0.0}
\newcommand{\FailedAfter}{0.0}
\newcommand{\VoiceMean}{0.4}
\newcommand{\VoiceMin}{0.3}
\newcommand{\VoiceMax}{1.5}
\newcommand{\UtilMax}{94.166}
\newcommand{\AltPathLatency}{0.0}
\newcommand{\RecoveryTrials}{5}
\newcommand{\RecoveryAfterMin}{25.3}
\newcommand{\RecoveryAfterMean}{26.7}
\newcommand{\RecoveryAfterMax}{31.9}
\newcommand{\RecoveryClockMin}{5.75}
\newcommand{\RecoveryClockMean}{5.85}
\newcommand{\RecoveryClockMax}{6.15}
\newcommand{\PartialTrials}{5}
\newcommand{\PartialAfterMin}{25.3}
\newcommand{\PartialAfterMean}{25.4}
\newcommand{\PartialAfterMax}{25.4}
\newcommand{\PartialClockMean}{5.92}
\newcommand{\OfferedTrials}{5}
\newcommand{\OfferedAfterMin}{25.0}
\newcommand{\OfferedAfterMean}{25.2}
\newcommand{\OfferedAfterMax}{25.3}
\newcommand{\OfferedClockMean}{6.22}
\newcommand{\NoOracleTrials}{5}
\newcommand{\NoOracleAfterMin}{25.3}
\newcommand{\NoOracleAfterMean}{25.3}
\newcommand{\NoOracleAfterMax}{25.5}
\newcommand{\NoOracleClockMean}{5.78}
\newcommand{\BaselineSamplePeak}{46.0}
\newcommand{\CongestedSamplePeak}{193.0}
\newcommand{\SevPhaseMean}{97.8}
\newcommand{\SevBeforeMean}{97.7}
\newcommand{\AlwaysActLatency}{25.7}
\newcommand{\AlwaysActBase}{17.6}
\newcommand{\AlwaysActPct}{+46}

% Numbers for the LLM-active replication (VIII-E).
\newcommand{\LlmModel}{qwen/qwen3.8-27b}
\newcommand{\LlmParseOk}{23}
\newcommand{\LlmNarrationRuns}{21}
\newcommand{\LlmDiagnoseRuns}{21}
\newcommand{\LlmParsedByModel}{23}
\newcommand{\LlmParseFallback}{0}
\newcommand{\LlmFallbackNarrations}{0}
\newcommand{\LlmRuns}{25}
\newcommand{\LlmOutcomesMatch}{25}
\newcommand{\LlmClockMean}{7.1}
\newcommand{\LlmClockMax}{12.2}

% Numbers for the boundary probe (Section VIII-B).
\newcommand{\ProbeTotal}{20}
\newcommand{\ProbeApproved}{9}
\newcommand{\ProbeRejected}{11}
\newcommand{\ProbeTelN}{8}
\newcommand{\ProbeTelApproved}{8}
\newcommand{\ProbeTelCross}{0}
\newcommand{\ProbeTopoN}{8}
\newcommand{\ProbeTopoApproved}{1}
\newcommand{\ProbeTopoCross}{6}
\newcommand{\ProbeTopoFaultyApproved}{0}
\newcommand{\ProbeCraftedN}{4}
\newcommand{\ProbeCraftedApproved}{0}
\newcommand{\ProbeCraftedCross}{1}
\newcommand{\ProbeStrong}{7}
\newcommand{\ProbeStrongApproved}{0}
\newcommand{\PromptCases}{6}


\begin{document}

\title{Trustworthy Agentic Automation for Intent-Based Self-Healing}
\author{\begin{tabular}{@{}c@{\hspace{3.2em}}c@{\hspace{3.2em}}c@{\hspace{3.2em}}c@{}}
Kamal Nayan Kumar & Altaf Raja & Ankit Kumar Keshari & Soumya Nimje \\[1pt]
{\small 23BDS026} & {\small 23BCS011} & {\small 23BCS015} & {\small 23BEC054} \\[5pt]
\multicolumn{4}{@{}c@{}}{Instructor: Dr. Jagadeesha R B}
\end{tabular}}
\maketitle

\begin{abstract}
Software-Defined Networking gives an operator precise programmatic control
over traffic paths, but it does not say \emph{when} to use that control, nor
whether a change helped. We present a locally-running system in which an
administrator states a performance goal in ordinary language and an agent
either restores the network, explains why it declined to act, or reports
that it could not measure honestly---every outcome grounded in a live
OpenFlow measurement, not an assertion.

The decision to change the network is never made by a probabilistic model:
a deterministic graph planner proposes, and three deterministic gates stand
between the candidate action and the data plane---an \emph{evidence} gate
that refuses missing, stale or insane measurements, a \emph{safety} gate
validating every mutation against nine constraints, and a \emph{verification}
gate re-measuring the network afterwards. A successful command is never
reported as a successful recovery. The system runs entirely on one machine
(Mininet, four OpenFlow 1.3 switches, Ryu controller, pure-Python datapath).
Across eight scenarios the agent produces the correct outcome for its
condition: a confident no-op, a verified recovery
(\BaselineMean{}~ms$\to$\CongestedMean{}~ms, repaired to \RecoveryAfter{}~ms),
an honest \emph{partial improvement}, a refusal to reroute a healthy flow,
a refusal of a request disabling a safety gate, an explicit report that the
network could not be measured, a verified recovery on a $60$~Mbit/s-configured
bottleneck crossed by real UDP load past $80\%$ utilisation, and the correct
localisation of an unflagged degraded link. We report three ablations and a
replication with the model in the loop (Section~VIII-B). The gates are
identical in both configurations; we state plainly what cannot yet be shown:
no adversarial models, a userspace Python switch, and lab restarts between
trials.
\end{abstract}

\begin{IEEEkeywords}
Software-Defined Networking, intent-based networking, network automation,
large language models, trustworthy AI, self-healing networks, verification.
\end{IEEEkeywords}

% ===========================================================================
\section{Introduction}
Software-Defined Networking (SDN) separates routing decisions (control
plane) from packet forwarding (data plane); a central controller programs
forwarding tables via OpenFlow, so paths change in milliseconds network-wide
\cite{mckeown2008openflow, feamster2014road, kreutz2015sdn}. That power is
incomplete without a second half: knowing \emph{when} rerouting is correct,
and whether it helped. A well-meaning reroute can degrade traffic or break
connectivity---and be reported as executed successfully.

LLMs are suited to the fuzzy parts: reading ``optimise video latency below
$50$~ms'' into a structured goal, and explaining a diagnosis. They cannot be
trusted with decisions that affect a running network, because they may
hallucinate a measurement, claim a path is safe without checking, or state
that a change succeeded unverified \cite{amodei2016safety}. We build an agent
that is useful \emph{and} bounded: \emph{the planner may propose; only
verified evidence may decide}. The pipeline is detection, evidence gathering,
decision (abstain if unjustified), path selection, safety validation,
execution, and verification---and steps 3, 5, 7 are the contribution.

Scope, before any claim: (i) candidate reroutes come from the deterministic
planner (the model, when configured, contributes intent and prose, never a
path); (ii) the headline evaluation has no model endpoint configured---the
agent parses intent deterministically---so the headline numbers are
reproducible without credentials. A replication with the model active
(\texttt{qwen/qwen3.8-27b} on Groq) is reported in Section~VIII-B; the
gates, planner and guardrail are identical in both configurations.

Our contributions are: (i) a runnable artifact whose three gates are plain
deterministic Python, reproducible and auditable; (ii) a live OpenFlow
evaluation measuring outcomes, including honest failures and refusals;
(iii) three ablations over the same recorded data; (iv) a truth audit
correcting twenty factual claims and fixing twenty-seven code defects
(\Repo AUDIT.md).

\section{Related Work}

\subsection{SDN and intent-based networking}

OpenFlow \cite{mckeown2008openflow} made programmable forwarding practical;
standard architecture surveys catalogue the design space and failure modes
\cite{feamster2014road,kreutz2015sdn}. Intent-based networking (IBN)
reframes configuration as outcome rather than command, reconciled against
observed state \cite{leivadeas2023ibn,rfc9315}. IBN surveys leave open what
a system may do when the objective is not met; the three gates are a
concrete answer, and the refusal is as much a part of it as the repair.

\subsection{Verification and safe control}

VeriFlow \cite{khurshid2013veriflow} and NetPlumber \cite{kazemian2013netplumber}
check \emph{static} configuration invariants pre-installation. Our safety
gate is of that kind; our verification gate is not---it re-measures the
network \emph{after} the change. Prior closed-loop systems divide
accordingly: ONOS's Intent Monitor and Reroute \cite{sanvito2018onos}
re-installs an intent on an alternate path when a constraint is violated
(no post-mutation outcome verdict, no refusal taxonomy); Lumi compiles
natural-language intents through Nile and requires operator confirmation
before deployment \cite{jacobs2021lumi}; NetConfEval \cite{wang2024netconfeval}
benchmarks LLM configuration synthesis; Confucius \cite{wang2025confucius}
and Protogeros et al.\ \cite{protogeros2026repair} place a model in the loop
with external checkers; STRATUS \cite{chen2025stratus} organises agents in a
state machine with an explicit \emph{Transactional No-Regression} safety
specification---the closest published statement of what our gates approximate,
and the rollback discipline we lack (Section~X). Our contribution is an
untrusted planner behind a deterministic runtime boundary that measures
\emph{afterwards}, refuses, abstains, and records why.

\subsection{LLMs for network operation}

Verified Prompt Programming \cite{mondal2023verified} pairs an LLM with
verifiers and localised feedback---the model is useful only because something
else checks it. The same shape appears at larger scale in NetConfEval
\cite{wang2024netconfeval}, \cite{wang2025confucius} and
\cite{protogeros2026repair}, and in NetIntent
\cite{hossain2025netintent}, which carries an LLM across the whole
intent-to-configuration path. The novelty is not a model that proposes
without checking, but a bounded pipeline where refusal is a first-class,
audited outcome.

\section{Problem Scenario}

Video between two sites is promised under a $50$~ms budget. An unrelated
background transfer crosses the video's path, the link congests, and
latency rises. Every operator decision in that state is made under
uncertainty with no automatic check that the guess was correct---a reroute
that appears to succeed but pushes latency from $97$~ms to $140$~ms looks
identical to one that repaired it, unless something measures afterwards.
The mirror failure is not acting: a single sub-threshold reading hides a
flow oscillating above and below the limit. The third failure is acting on
a healthy network because something \emph{wants} action. This paper
demonstrates against all three: a four-host, four-switch Mininet lab with
Ryu and a pure-Python OpenFlow 1.3 datapath, driven by a FastAPI backend
and a LangGraph agent (Fig.~\ref{fig:topology}). The congestion symptom on
$s2$-$s4$ is emulated by injecting link delay (most scenarios) or by real
background UDP (S7); $h2$ carries that traffic to $h4$.

\section{System Model}

\subsection{Network model}
The network is a graph $G=(V,E)$, $V=H\cup S$ (hosts, switches); each edge
has a one-way delay $d_e$, a capacity $c_e$, and a status. \emph{Switch
edges} $E_S$ are the only telemetry-bearing ones; host access edges add a
fixed $0.1$~ms. A \emph{flow} is $(s,t)\in H\times H$. A \emph{path} is
$p=(v_0,\dots,v_k)$, $v_0=s$, $v_k=t$, with switch-edge set
$E_p=\bigl\{\{v_i,v_{i+1}\}:0\le i<k\bigr\}\cap E_S$ and one-way
switch delay
\begin{equation}
D(p)=\sum_{e\in E_p} d^{\mathrm{eff}}_e ,
\label{eq:delay}
\end{equation}
where $d^{\mathrm{eff}}_e$ is the delay the link currently carries
(Section~IV-B): the configured value normally, the larger in-force value
while a scenario holds the link. The ICMP round trip is
\begin{equation}
\mathrm{RTT}(p)\;\approx\;2\,\bigl(D(p)+2\,\delta_{\mathrm{acc}}\bigr)
+c_{\mathrm{dp}},
\label{eq:rtt}
\end{equation}
with $\delta_{\mathrm{acc}}=0.1$~ms. $c_{\mathrm{dp}}$ is \emph{not}
assumed: the calibration run measured $17.3$~ms ($c_{\mathrm{dp}}\approx
0.9$~ms), this submission's run $\BaselineMean$~ms ($\approx1.2$~ms)---the
same order, and the reason every claim compares against measured values,
never Eq.~\eqref{eq:rtt}. An earlier datapath never matched a rule and gave
$c_{\mathrm{dp}}\approx 13.7$~ms (Section~X).

The controlled intent is a tuple
\begin{equation}
I=(s,t,\tau,\ell,\mathcal{P}),
\label{eq:intent}
\end{equation}
where $\tau>0$ is the latency target in ms, $\ell$ an optional loss target,
and $\mathcal{P}$ the protected flows; an element $q=(s_q,t_q,\tau_q,
\ell_q,p_q)$ gives its endpoints, budgets, and path. The evaluated intent
declares one protected flow, the voice flow $q=(\texttt{h3},\texttt{h4},
20\,\mathrm{ms},1\,\%,[\texttt{h3},\texttt{s4},\texttt{h4}])$.

\subsection{Observation model}
At any time the lab produces a \emph{snapshot}
$\Sigma=(G,\mathcal{L},\mathcal{F},C_{\mathrm{deg}})$: the topology $G$;
per-link telemetry $\mathcal{L}=\{\ell_e\}$ with
$\ell_e=(u_e,c_e,d^{\mathrm{eff}}_e,\mathrm{status}_e)$; per-flow
measurements $\mathcal{F}=\{f\}$ (endpoints, path, sample distribution
$\mathbf{s}_f$, peak, loss); and the set $C_{\mathrm{deg}}$ of degraded or
no-longer-up links. $C_{\mathrm{deg}}$ and the gate's suspect set
$C_{\mathrm{susp}}$ (Section~VI-B) are kept apart: the former is a lab fact
used by Eq.~\eqref{eq:score} and the capacity check; the latter is the
gate's \emph{diagnosis}, used by the planner. S8 removes the flag; S2/S3/S7
coincide. Every element is measured: latency by in-host ICMP, utilisation by
the switch's port counters, status by the controller. $d^{\mathrm{eff}}_e$
is a per-link value \emph{reported by the lab} (configured delay plus
injected delay), not inferred from host probes; Section~X and roadmap
item~4 state what that costs us.

\subsection{Trust boundary}
The agent is a state graph over typed state. $M$ (the LLM) produces text and
structured hints; deterministic components $D$ (gates, planner, executor)
contain no model call. The only edge from $M$ into $D$ is the typed intent,
validated against the schema and topology; the only edge from $D$ into the
data plane passes through the safety gate.

\section{Mathematical Formulation}
% ===========================================================================

\subsection{Violation}

Let $f$ be the measured flow for intent $I$, with the non-negative ping
samples $\mathbf{s}_f=(s_1,\dots,s_n)$, its reported peak, and its loss
$p_f$. We define the violation predicate with three outcomes---violated,
not violated, and \emph{unknown}:

\begin{equation}
v(f,I)=
\begin{cases}
\text{unknown}, & \text{if } f \text{ is absent or unmeasured,}\\
\text{true},     & \text{if } n \ge 2,\ \bigl|\{i:s_i>\tau\}\bigr| \ge \gamma n,\\
\text{false},    & \text{if } n \ge 2,\ \bigl|\{i:s_i>\tau\}\bigr| < \gamma n,\\
\text{true},     & \text{if } n < 2 \text{ and the peak } > \tau,\\
\text{true},     & \text{if } n < 2,\ \text{no peak, latency } > \tau,\\
\text{false},    & \text{otherwise},
\end{cases}
\label{eq:violation}
\end{equation}
with $\gamma=0.25$ the fraction of samples that may exceed the target,
lines evaluated in order. Packet loss is deliberately \emph{absent}: an
intent may declare an optional loss target $\ell$, but the violation
predicate is a statement about the latency distribution only; loss is
evaluated in the verification verdict (Section~V-D via $L$), so the two
predicates have disjoint scopes and the third row of Eq.~\eqref{eq:verify}
is reachable. The second and third lines are the operative ones: an SLO is
a statement about the \emph{distribution} of delays, not their average;
$\bar{s}$ alone makes the verdict a coin flip on a flow oscillating around
the target (Section~IX-C). The last three lines handle $n<2$: one sample
falls back to the reported peak, then to the reported latency (for $n=1$,
that single sample); two samples already face the distribution rule. The
three-way outcome is what lets the system distinguish a confident no-op
from a refusal.

Gate~1's verdict $e(\Sigma,I)$ is likewise a conjunction: topology present
and both endpoints in it; the flow measured (latency and loss present);
page $\le 5$~s; samples $0\le s_i\le 10{,}000$~ms, $0\le p_f\le 100\%$,
no NaN; at least one named suspect link present in live telemetry and
topology; and confidence $\ge 0.6$, where confidence $=
0.6\cdot(\mathrm{present})+0.4\cdot(\mathrm{fresh})$ over evidence items.
Failures yield \texttt{insufficient}, \texttt{stale} or
\texttt{unavailable}; the floor and freshness are configuration, not code.

\subsection{Candidate ranking}

$\mathcal{P}_c$ is the candidate set for flow $(s,t)$. A candidate is
\emph{congestion-free} when it avoids both the lab's degraded set and the
gate's suspects:
\begin{equation}
A(p)=\bigl[E_p\cap (C_{\mathrm{deg}}\cup C_{\mathrm{susp}})=\varnothing\bigr].
\label{eq:avoids}
\end{equation}
Only a strictly-better candidate is worth installing, so
\begin{equation}
\mathrm{improve}(p) \;\equiv\; D(p) < D(p_{\mathrm{cur}}),
\label{eq:improve}
\end{equation}
(a pre-filter, not a safety check: evaluated first, tagged
\texttt{not\_an\_improvement}. Without it a planner ranking only by
Eq.~\eqref{eq:score} would install a slower path over a healthy one ---
scenario S4).

We score candidates so that lower is better:
\begin{equation}
\mathrm{score}(p)=
\begin{cases}
+\infty, & \neg A(p),\\[2pt]
\displaystyle\sum_{e\in E_p} d^{\mathrm{eff}}_e
 + \frac{1}{100}\max_{e\in E_p} u_e, & A(p).
\end{cases}
\label{eq:score}
\end{equation}
The sum uses the \emph{effective} delay (so a degraded link is expensive
even with zero byte counters), runs over \emph{all} of $E_p$, and the
disqualifying set includes the gate's suspects---so an unnamed bottleneck
cannot silently strip congestion-awareness. S8 removes the lab flag; the
repair still succeeds (Section~VIII). Section~X states the remaining
dependency.
Eq.~\eqref{eq:avoids} is a filter, not just a score: the planner's
enumeration already excludes $C_{\mathrm{deg}}\cup C_{\mathrm{susp}}$, so
a candidate failing $A(p)$ never reaches Gate~2; Eq.~\eqref{eq:score}'s
$+\infty$ is the score such a path would carry if it were scored.
Consequently the nine checks only ever see suspect-free candidates (Check~6
can never be the check that excludes a suspect edge---in any scenario),
and the gate's independent value lies in what the planner cannot pre-empt
(endpoints, path shape, hop budget, flow scope), measured in Section~VIII-C.

Crossing a degraded link is \emph{disqualifying} rather than merely
expensive: routing a flow around a bottleneck is the entire point, so a path
that still crosses the bottleneck is not a repair. Among clean paths the
accumulated delay decides and utilisation only breaks ties. Section~IX-B
measures what this rule is worth against the obvious alternative of ranking
on delay alone.

\subsection{Safety validation}

A proposed action $a=(\mathrm{op}, s, t, p)$ is admitted only if all nine
checks pass ($\mathrm{op}(a)\in\{\texttt{reroute\_flow}\}$, the only op
the system ever proposes). Checks:

\begin{enumerate}
\item \textbf{operation\_allowlisted}: $\mathrm{op}(a)\in
      \{\texttt{reroute\_flow}\}$, the complete allow-list;
\item \textbf{endpoints\_exist}: $s,t\in V$;
\item \textbf{path\_well\_formed}: $p=(v_0,\dots,v_k)$ with $v_0=s$,
      $v_k=t$, no repeated node, and every node of $p$ in $V$;
\item \textbf{hop\_budget}: the number of switch hops $k-1\le h_{\max}$,
      with $h_{\max}=6$;
\item \textbf{all\_links\_operational}: every $e\in E_p$ exists in
      $E(G)$ and is up, both in the recorded topology and in the live
      telemetry;
\item \textbf{capacity\_respected}: no $e\in E_p$ is degraded
      ($e\notin C_{\mathrm{deg}}$) and none has $u_e$ at
      or above the congestion threshold, which is $80\%$ of capacity;
\item \textbf{protected\_flows\_unmodified}: $E_p\cap
      E_{\mathrm{prot}}=\varnothing$;
\item \textbf{flow\_scope}: the action's endpoints are the \emph{intent's}
      endpoints, not merely the path's;
\item \textbf{request\_not\_expired}: $\mathrm{age}(a)\le 300\,\mathrm{s}$.
\end{enumerate}

Check 6 treats a degraded link as a capacity violation even at zero
utilisation---the injected fault is delay, invisible to counters. Check 7
compares the inter-switch edges, not nodes (two flows sharing a destination
share its access switch, so a node test would reject every repair). Check 8
distinguishes action from intent: a well-formed reroute of the \emph{wrong
flow} is refused.

Check 6 tests only $e\notin C_{\mathrm{deg}}$ and $u_e<80\%$---it neither
reads $C_{\mathrm{susp}}$ nor enforces an effective-delay bound. The
predicate that would close it is simple and we state it even though the
implementation does not enforce it: $e\in E_p\Rightarrow
d^{\mathrm{eff}}_e\le\alpha d_e$ (proposed $\alpha=2$), equivalently
$E_p\cap C_{\mathrm{susp}}=\varnothing$. In S8, $C_{\mathrm{deg}}=
\varnothing$ and $u_e=0\%$, so Check~6 passes vacuously for every
candidate, including $s1$-$s2$-$s4$ across the fault; S8 succeeds only
because Eq.~\eqref{eq:avoids} removed that path upstream. Gate~2 is the
planner's backstop, not an independent opinion on unflagged faults
(Section~VIII-C quantifies the gap; roadmap item~4 closes it).

The checks are pure functions of a snapshot and the action---no model call,
deterministic verdicts, each named check recorded separately (an
unattributable rejection is not auditable). Risk tiers are assigned from
the checks: HIGH (protected-flow or allowlist failure), MEDIUM (capacity,
path-well-formedness, flow-scope), otherwise LOW. A \emph{failed} check
rejects; the hold path (\texttt{awaiting\_approval}) applies only to an
\emph{approved} action whose tier is in the unattended-hold set (default
MEDIUM/HIGH). In the evaluated configuration, approved actions always pass
everything, so the hold path is never taken---stated rather than hidden.

\subsection{Verification}

Let $f^-,f^+$ be the flow measurements before and after the action,
$b,a$ their reported (scalar) mean RTTs, $p_{f^+}$ the post-action loss,
$L\equiv(\ell\text{ undeclared})\vee(p_{f^+}\le\ell)$ the loss target.
$\mathcal{Q}$ is the (possibly empty) set of protected-flow evaluations;
$Q$ is their conjunction, vacuously true when empty. A protected flow $q$
is ok when its measured latency is within its constraint, its loss within
its limit, and its path unchanged. $v(f^+,I)$ is Eq.~\eqref{eq:violation}
on the post-action measurement ($\neg v$ reads the latency branches; $L$
carries the loss constraint). The verdict applies the first matching rule:

\begin{equation}
\mathcal{V}=
\begin{cases}
\texttt{unverified}, & \text{if } f^+ \text{ absent or unmeasurable},\\[2pt]
\texttt{success},    & \text{if } \neg v(f^+,I) \wedge L \wedge Q,\\[2pt]
\texttt{partial},    & \text{if } \neg v(f^+,I) \wedge \neg L \wedge Q,\\[2pt]
\texttt{failed},     & \text{if } \neg Q,\\[2pt]
\texttt{partial},    & \text{if } b \text{ known and } a < b,\\[2pt]
\texttt{failed},     & \text{otherwise}.
\end{cases}
\label{eq:verify}
\end{equation}
where the conditions in each row are conjunctions, so \texttt{failed} fires
whenever a protected flow is outside its constraints, whatever happened to
the latency.

Three ordering decisions are deliberate. Row~3 (latency within target, loss
not) is distinct from a degenerate row~5: the honest report is \emph{partial},
not success. Row~5 (improvement) fires only after row~4 (protected flows), and
on \emph{improvement} regardless of whether the target is met---reporting a
partial as success is the most damaging lie such a system could tell. Row~1 is
not a verdict at all, and the implementation keeps it apart from
``it did not recover''.

The 1-to-1 correspondence between Eq.~\eqref{eq:verify} and the terminal
outcomes the implementation records is:
\begin{center}
\footnotesize
\begin{tabular}{ll}
\toprule
Eq.~\eqref{eq:verify} & recorded outcome \\
\midrule
\texttt{unverified} & \texttt{unverified} \\
\texttt{success}    & \texttt{intent\_satisfied} \\
\texttt{partial}    & \texttt{partial\_improvement} \\
\texttt{failed}     & \texttt{recovery\_failed} \\
\bottomrule
\end{tabular}
\end{center}
No evaluated run produced \texttt{unverified}; it exists in the taxonomy and
is recorded distinctly rather than folded into \texttt{recovery\_failed}.

\emph{The gate judges post-action latency by the same
sample-distribution predicate Eq.~\eqref{eq:violation} uses, not a scalar
mean.}

% ===========================================================================
\section{Proposed Algorithm}
% ===========================================================================

\subsection{Overview}

An explicit LangGraph state graph (Fig.~\ref{fig:workflow}) with thirteen
compiled nodes and conditional routing: \texttt{intent}, \texttt{observe},
\texttt{evidence}, \texttt{diagnose}, \texttt{validate}, \texttt{execute},
\texttt{observe\_after}, \texttt{verify}, plus terminals \texttt{abstain},
\texttt{satisfied}, \texttt{stop}, \texttt{hold}, \texttt{finish}. Network
mutation is fixed in code; the model is consulted only where errors are
recoverable. The routing is the implementation's routing function, not a
paraphrase (Algorithm~\ref{alg:agent}, p.~\pageref{alg:agent}). The red
boxes in Fig.~\ref{fig:workflow} are the five compiled terminals, drawn by
\emph{outcome} rather than node: \texttt{stop} carries
\texttt{guardrail\_blocked}, \texttt{clarification}, \texttt{rejected},
\texttt{no\_better\_path}; \texttt{finish} is post-verify only; early
exits skip it. The evidence gate routes to \textbf{satisfied} iff the
violation verdict is false, the loss predicate holds, and evidence is
sufficient/fresh/sane; to \textbf{diagnose} iff evidence is sufficient and
the intent is violated; else \textbf{abstain}. The safety gate routes to
\textbf{execute} iff approved and not approval-gated, to \textbf{hold} iff
approved but approval-gated, to \textbf{stop} otherwise.

A failed \texttt{observe} sends the run to the \texttt{abstain} terminal
but preserves the \texttt{lab\_offline} outcome---``the network is fine''
and ``the network could not be read'' are not the same report.

\begin{figure}[t]
  \centering
  \includegraphics[width=\columnwidth]{figures/fig_workflow.pdf}
  \caption{The state graph. Amber nodes are the deterministic gates.
  The thirteen compiled nodes are drawn with three \emph{outcome}-labelled
  terminal boxes: \texttt{stop} carries \texttt{guardrail\_blocked},
  \texttt{clarification} and \texttt{rejected}; \texttt{abstain} is the
  evidence-short terminal; \texttt{satisfied} is \texttt{intent\_satisfied};
  \texttt{hold} is \texttt{awaiting\_approval}; \texttt{finish} is reached
  only after verification. Early exits bypass it.}
  \label{fig:workflow}
\end{figure}

\begin{algorithm}[t]
\small
\caption{The guarded recovery workflow. The node names are those of the
compiled graph, and the routing is the routing function of the
implementation, not a paraphrase of it. Every branch terminates in exactly
one terminal node; the API persists the run record from the terminal state.}
\label{alg:agent}
\begin{algorithmic}[1]
\Require snapshot function $\Sigma(\cdot)$, validator $V$, executor $X$
\Ensure a terminal outcome $o$: \texttt{intent\_satisfied},
  \texttt{partial\_improvement}, \texttt{recovery\_failed},
  \texttt{unverified}, \texttt{no\_better\_path}, \texttt{action\_rejected},
  \texttt{awaiting\_approval}, \texttt{abstained\_insufficient\_evidence},
  \texttt{lab\_offline}, \texttt{guardrail\_blocked}, or a clarification
\State \textbf{intent}: screen the request with the deterministic guardrail
  and parse $I$ (model, else deterministic fallback)
\If{the request tries to override the instructions, move protected traffic,
    or disable a gate}
  \State $o\gets\texttt{guardrail\_blocked}$ via \texttt{stop}; \Return
  \Comment{no telemetry is read}
\EndIf
\If{$I$ is ambiguous}
  \State ask the operator to clarify; $o\gets$ clarification via
    \texttt{stop}; \Return
\EndIf
\State \textbf{observe}: $\Sigma_0\gets\Sigma(I)$
\If{$\Sigma(I)$ fails (the lab is unreachable)}
  \State $o\gets\texttt{lab\_offline}$ via \texttt{abstain}; \Return
  \Comment{the reason survives: abstention does not overwrite it}
\EndIf
\State \textbf{evidence}: $v\gets v(f,I)$ (Eq.~\eqref{eq:violation}),
  $L\gets$ the loss predicate of Section~V-D, $(e,\text{freshness})\gets$
  the evidence verdict
\If{$v=\text{false} \wedge L \wedge e=\texttt{sufficient}$}
  \State $o\gets\texttt{intent\_satisfied}$ via \texttt{satisfied}; \Return
  \Comment{confident no-op; freshness/sanity checked first}
\ElsIf{$e=\texttt{sufficient}$}
  \State go to \textbf{diagnose}
\Else
  \State $o\gets\texttt{abstained\_insufficient\_evidence}$ via
    \texttt{abstain}; \Return \Comment{no plan, no change}
\EndIf
\State \textbf{diagnose}: $C_{\mathrm{susp}}\gets$ the suspect links named
  by Gate~1 (congested links first, else all on-path links, slowest
  effective delay first), $\mathcal{P}_c\gets$ candidate paths avoiding
  $C_{\mathrm{deg}}\cup C_{\mathrm{susp}}$, ranked by Eq.~\eqref{eq:score};
  drop any candidate equal to the path the flow is already on
\State \textbf{validate}: \emph{pre-filter} $\mathcal{P}_c$ by the
  improvement test---a candidate whose effective switch delay is not less
  than the current path's is dropped here, before any safety check is
  evaluated; every drop is recorded with
  \texttt{failed\_checks = [not\_an\_improvement]}
\If{$\mathcal{P}_c=\varnothing$ on entry (the topology offers no
    alternative path)}
  \State $o\gets\texttt{action\_rejected}$ via \texttt{stop}; \Return
  \Comment{no usable alternative path exists}
\EndIf
\If{no candidate survives the pre-filter and every recorded rejection is a
    non-improvement}
  \State $o\gets\texttt{no\_better\_path}$ via \texttt{stop}; \Return
  \Comment{declined: rerouting would make the flow worse}
\EndIf
\State take the first $p\in\mathcal{P}_c$ that $V$ admits (the nine checks
  of Section~V-C); record every rejection
\If{no $p$ is admitted}
  \State $o\gets\texttt{action\_rejected}$ via \texttt{stop}; \Return
\ElsIf{$V$ requires human approval}
  \State $o\gets\texttt{awaiting\_approval}$ via \texttt{hold}; \Return
  \Comment{nothing changed}
\Else
  \State \textbf{execute}: $X(p)$ \Comment{the only mutation}
  \State \textbf{observe\_after}: $\Sigma_1\gets\Sigma(I)$ \Comment{fresh}
  \State \textbf{verify}: $o\gets\mathcal{V}(I,f^-,f^+,Q)$
    (Eq.~\eqref{eq:verify}) \Comment{Eq. (7) rows map to the outcomes of Section~V-D}
\EndIf
\State \texttt{finish} (post-verification only) stamps the finish time;
  the API persists the run record from the terminal state; \Return $o$
\end{algorithmic}
\end{algorithm}

\subsection{Gate 1: evidence sufficiency}

Gate~1 checks the measurements exist, are recent ($\le 5$~s), are sane, and
support diagnosis. A stale snapshot satisfying $\tau$ is routed to
abstention, never to a confident no-op. The suspect rule: with no congested
link ($u_e\ge 80\%$), name the congested links (most-utilised first);
otherwise name \emph{all} current-path links by effective delay, slowest
first. In S2 the degraded path $s1$-$s2$-$s4$ has
$d^{\mathrm{eff}}_{s2\text{-}s4}=4+40=44$~ms and $s1$-$s2$ $4$~ms (set
$\{s2\text{-}s4,s1\text{-}s2\}$, slowest first; one-way delay
$D(s1\text{-}s2\text{-}s4)=48$~ms). Including the healthy upstream link
disqualifies otherwise-clean candidates sharing it ($s1$-$s2$-$s3$-$s4$)---
the cost of a conservative diagnosis, stated, not hidden; S8 depends on the
same mechanism and succeeds. Lab delays: $s1$-$s2$ $4$, $s2$-$s4$ $4$
(100~Mbit/s, $60$ in S7), $s1$-$s3$ $5$, $s2$-$s3$ $4$, $s3$-$s4$ $7$~ms.

\subsection{Gate 2: deterministic action safety}

Gate~2 returns a verdict, a risk tier, and the named checks. It never
consults the model. A degraded link fails capacity even at zero utilisation.
Scenario S8 shows the limit plainly: Check~6 passes vacuously on unflagged
delay faults ($e$ not in $C_{\mathrm{deg}}$, $u_e$ below threshold); the
path across the fault is absent only because the planner excluded it
upstream. Gate~2 is a backstop, not an independent opinion on unflagged
faults.

\subsection{Gate 3: recovery verification}

The verification gate re-measures the flow and applies Eq.~\eqref{eq:verify}.
It distinguishes success, partial, failed, unverified---and treats a missing
measurement as unverified, not success. Protected flows are checked
\emph{before} the improvement test: harming one is a failure, however much
the intent flow improved.

\subsection{Deterministic intent parsing and the guardrail}

With no model configured, a keyword parser extracts the stated endpoints,
target and traffic class only; unstated facts keep documented defaults. A
deterministic guardrail screens the request first: no identity overrides, no
instruction overrides, no gate disables, no protected-traffic moves---a
refusal is published as loudly as an action.

% ===========================================================================
\section{Baselines and Experimental Setup}
% ===========================================================================

\subsection{The lab}

The lab is live: four hosts, four OpenFlow 1.3 switches in Mininet,
Ryu controller, unprivileged network namespace; the datapath is the project's
pure-Python switch (Open vSwitch unavailable; the backend reports which one
is in use, so nothing is silently substituted). The lab API lives on a unix
domain socket; no TCP port is opened. $s2$-$s4$ is the bottleneck (default
$100$~Mbit/s; S7 reconfigures it to $60$~Mbit/s so UDP load crosses the
$80\%$ utilisation threshold on live byte counters). Configured delays:
$s1$-$s2$ $4$, $s2$-$s4$ $4$, $s1$-$s3$ $5$, $s2$-$s3$ $4$, $s3$-$s4$
$7$~ms; access edges $0.1$~ms each. The four simple paths: $8$~ms,
$12$~ms, $15$~ms, $13$~ms. An earlier datapath revision never matched an
installed rule (every frame forwarded via the controller, $c_{\mathrm{dp}}
\approx 13.7$~ms); it is fixed and regression-tested ($c_{\mathrm{dp}}
\approx 0.9$~ms measured).

\subsection{The scenarios and the faults}

The fault is synthetic, not organic: $40$~ms of \texttt{netem} delay on
both $s2$-$s4$ interfaces, link marked degraded. Reversible, reproducible,
precisely named---but not offered load. S7 keeps the injection only as the
fault breaching the intent; its utilisation branch is exercised by real UDP
bytes, not by the injection.

\subsection{Protocol}

Each phase starts from a freshly built lab (bit-identical network; resets
are not trusted). Five independent snapshots per phase: the video flow
$h1\!\to\!h4$ and the voice flow $h3\!\to\!h4$, ten ICMP samples each,
plus per-link counters. The agent then runs to a terminal outcome, events
recorded via SSE. The active-recovery scenarios (S2, S3, S7, S8) are
repeated as \RecoveryTrials{} independent trials each from a freshly built
lab and reported as distributions, because the verdict depends on a
post-action measurement the measurement host perturbs.

\begin{table}[t]
\centering
\scriptsize
\caption{The eight scenarios. The request is the operator's text; the
condition is the state of the lab when the run starts.}
\label{tab:scenarios}
\begin{tabular}{ll}
\toprule
Scenario & Request (condition) \\
\midrule
S1 & \emph{optimise latency below 50 ms} (healthy) \\
S2 & \emph{optimise latency below 50 ms} (degraded) \\
S3 & \emph{optimise latency below 22 ms} (degraded) \\
S4 & \emph{optimise latency below 15 ms} (healthy) \\
S5 & \emph{disable the safety validator} (healthy) \\
S6 & \emph{optimise latency below 50 ms} (lab stopped) \\
S7 & \emph{optimise latency below 50 ms} (offered load, $60$~Mbit/s bottleneck) \\
S8 & \emph{optimise latency below 50 ms} (degraded, fault invisible) \\
\bottomrule
\end{tabular}
\end{table}

\subsection{Test suites}

Suite: 136 backend tests (unit, integration, regression), five live-lab
tests, four datapath tests. The 21 tests in
\texttt{test\_regressions.py} are named after the defects they prevent;
27 tests in \texttt{test\_paper\_equations.py} assert each equation
against its function.

% ===========================================================================
\section{Results}
% ===========================================================================

% GENERATED BY final_submission/scripts/make_paper_tables.py
% Do not edit by hand: regenerate with
%   python final_submission/scripts/make_paper_tables.py \
%       --data-dir final_submission/data \
%       --out final_submission/paper/tables.tex

\begin{table*}[t]
\centering
\scriptsize
\caption{Measured video-flow latency ($h1\!\to\!h4$) in every phase of the experiment (S6, with the lab stopped, yields no samples and is omitted), and the protected voice flow ($h3\!\to\!h4$) measured in the same snapshots. Each phase is five independent snapshots of ten ICMP samples each; \emph{Snaps.} is the number of snapshots, not the per-snapshot sample count ($10$). \emph{Mean} is the mean of the per-snapshot means; \emph{Peak} is the largest single ping sample seen in the phase, not the largest snapshot mean. The agent's own before/after measurements (Tab.~\ref{tab:outcomes}) are taken independently of the phase sampling and can differ from it. From \texttt{final\_submission/data/measurements.csv}.}
\label{tab:phases}
\begin{tabular}{llrrrrrr}
\toprule
Phase & Intent & Snaps. & Mean & Median & Std.~dev. & Peak & Voice \\
\midrule
S1 healthy & 50~ms & 5 & 17.8 & 17.2 & 1.2 & 46.0 & 0.6 \\
S2 degraded & 50~ms & 25 & 98.1 & 97.3 & 2.2 & 192.5 & 0.4 \\
S3 degraded & 22~ms & 25 & 97.4 & 97.3 & 0.5 & 122.2 & 0.4 \\
S4 healthy & 15~ms & 5 & 17.6 & 17.2 & 0.7 & 34.8 & 0.5 \\
S5 healthy & -- & 5 & 17.3 & 17.3 & 0.1 & 19.7 & 0.4 \\
S7 degraded+load & 50~ms & 25 & 97.8 & 97.7 & 0.5 & 111.3 & 0.4 \\
S8 degraded, no flag & 50~ms & 25 & 97.3 & 97.3 & 0.1 & 98.7 & 0.4 \\
S1' restart & -- & 5 & 17.3 & 17.3 & 0.1 & 18.0 & 0.4 \\
\bottomrule
\end{tabular}
\end{table*}

\begin{table*}[t]
\centering
\scriptsize
\caption{Outcome of every agent run. ``Before'' and ``after'' are the latencies the verification gate measured; ``--'' means the run never changed the network. The gate column is the number of the nine deterministic safety checks that passed. From \texttt{final\_submission/data/scenario\_summary.csv}.}
\label{tab:outcomes}
\begin{tabular}{llrrrl}
\toprule
Scenario & Outcome & Before & After & Installed path & Gates \\
\midrule
S1 healthy & \texttt{intent\_satisfied} & -- & -- & \texttt{--} & -- \\
S2 degraded & \texttt{intent\_satisfied (5 trials)} & 99.3 [97.2--107.4] & 26.7 [25.3--31.9] & \texttt{s1-s3-s4} & 9/9 \\
S3 degraded & \texttt{partial\_improvement (5 trials)} & 97.3 [97.2--97.4] & 25.4 [25.3--25.4] & \texttt{s1-s3-s4} & 9/9 \\
S4 healthy & \texttt{no\_better\_path} & -- & -- & \texttt{--} & -- \\
S5 guardrail & \texttt{guardrail\_blocked} & -- & -- & \texttt{--} & -- \\
S6 lab offline & \texttt{lab\_offline} & -- & -- & \texttt{--} & -- \\
S7 load & \texttt{intent\_satisfied (5 trials)} & 97.7 [97.0--98.3] & 25.2 [25.0--25.3] & \texttt{s1-s3-s4} & 9/9 \\
S8 degraded, no flag & \texttt{intent\_satisfied (5 trials)} & 97.3 [97.3--97.3] & 25.3 [25.3--25.5] & \texttt{s1-s3-s4} & 9/9 \\
\bottomrule
\end{tabular}
\end{table*}

\begin{table*}[t]
\scriptsize
\centering
\caption{Outcome distribution per scenario. The active-rerouting scenarios (S2 recovery, S3 partial improvement, S7 offered load, S8 no-oracle) are repeated as five independent trials, each from a freshly built lab; the other scenarios are run once. The latency column is the post-action latency the verification gate measured; the wall-clock column is the end-to-end closed-loop duration (pre-action snapshot to verified outcome), min--max over the trials. From \texttt{final\_submission/data/scenario\_aggregates.csv}.}
\label{tab:aggregates}
\begin{tabular}{llrrll}
\toprule
Scenario & Trials & Outcome(s) & After (min--max) & Wall-clock (min--max) & Path \\
\midrule
S1 healthy & 1 & \texttt{intent\_satisfied (1)} & -- & 1.9--1.9 s & \texttt{--} \\
S2 degraded & 5 & \texttt{intent\_satisfied (5)} & 25.3--31.9 ms & 5.7--6.2 s & \texttt{s1-s3-s4} \\
S3 degraded & 5 & \texttt{partial\_improvement (5)} & 25.3--25.4 ms & 5.8--6.2 s & \texttt{s1-s3-s4} \\
S8 degraded, no flag & 5 & \texttt{intent\_satisfied (5)} & 25.3--25.5 ms & 5.8--5.8 s & \texttt{s1-s3-s4} \\
S4 healthy & 1 & \texttt{no\_better\_path (1)} & -- & 3.8--3.8 s & \texttt{--} \\
S5 guardrail & 1 & \texttt{guardrail\_blocked (1)} & -- & 0.0--0.0 s & \texttt{--} \\
S6 lab offline & 1 & \texttt{lab\_offline (1)} & -- & 0.1--0.1 s & \texttt{--} \\
S7 load & 5 & \texttt{intent\_satisfied (5)} & 25.0--25.3 ms & 6.2--6.2 s & \texttt{s1-s3-s4} \\
\bottomrule
\end{tabular}
\end{table*}

\begin{table}[t]
\centering
\caption{Path selection under the degraded bottleneck. Both rows rank the \emph{same} candidate set the lab produced, so the comparison isolates the ranking rule. The measured latency is what the live network recorded on the selected path. From \texttt{final\_submission/data/ablations.csv}.}
\label{tab:ranking}
\begin{tabular}{llr}
\toprule
Ranking rule & Path & Latency \\
\midrule
system (congestion-aware) & \texttt{s1-s3-s4} & 25.4~ms \\
ablation (shortest path) & \texttt{s1-s2-s4} & 97.2~ms \\
\bottomrule
\end{tabular}
\end{table}

\begin{table}[t]
\centering
\caption{Recorded measurements on which the two violation rules disagree, one row per direction. The system treats a flow as violated when at least 25\% of its individual ping samples exceed the target; the mean-only ablation compares only the mean. \emph{System / mean-only} reports each rule's verdict on the intent (yes = violated). \texttt{yes / no} is the oscillation case the sample-distribution rule is built to catch; \texttt{no / yes} is the outlier-tolerance case (one high sample, below $\gamma$). Rows are the closest calls in the recorded data (smallest $|\bar{s}-\tau|$), from a millisecond sweep of $\tau$ from $5$ to $200$~ms, not an estimate of the deployed operating point. From \texttt{final\_submission/data/ablations.csv}.}
\label{tab:violation}
\tiny
\begin{tabular}{@{}l r r l >{\raggedright\arraybackslash}p{0.34\columnwidth}@{}}
\toprule
Phase & Target & Mean & Sys / mean & Samples (ms) \\
\midrule
S2 degraded & 98 & 98.1 & no / yes & {\ttfamily\scriptsize 105.0, 97.4, 97.2, 97.2, 97.4, 97.3, 97.3, 97.2, 97.3, 97.4} \\
S7 degraded+load & 98 & 98.0 & yes / no & {\ttfamily\scriptsize 96.8, 98.6, 97.6, 99.1, 99.0, 97.6, 98.4, 98.0, 97.9, 96.9} \\
\bottomrule
\end{tabular}
\end{table}

\begin{table}[t]
\centering
\caption{Acting when there is nothing to act on. The system ends a healthy run with no flow rules installed; the always-act variant takes the same code path the agent uses when a target is unreachable, so its cost is measured rather than asserted. From \texttt{final\_submission/data/ablations.csv}.}
\label{tab:alwaysact}
\scriptsize
\begin{tabular}{llr}
\toprule
Policy & Outcome & Healthy flow \\
\midrule
system & \texttt{no-op (nothing installed)} & 17.6~ms, unchanged \\
always act & \texttt{rerouted} & 17.6~$\to$~25.7~ms (+46\%) \\
\bottomrule
\end{tabular}
\end{table}



\subsection{Outcomes}

Tab.~\ref{tab:phases} gives the sampled measurements behind every phase, and
Tab.~\ref{tab:outcomes} gives the terminal outcome of every scenario. Each
outcome is the correct answer for its condition.

\textbf{S1.} \BaselineMean{}~ms (no-op, correct; peak
\BaselinePeak{}~ms). The agent returned \texttt{intent\_satisfied},
installed no rules, and said why.

\textbf{S2.} Latency rose to \CongestedMean{}~ms. The evidence gate named
$s2$-$s4$ (the scenario's link; slowest on path) and then $s1$-$s2$; Eq.~\eqref{eq:score} ranked the
degraded link last; the planner installed $s1$-$s3$-$s4$; all nine checks
passed; Gate~3 measured \RecoveryAfterMean{}~ms ($\texttt{intent\_satisfied}$).
Repeated as \RecoveryTrials{} fresh-lab trials:
\RecoveryAfterMin{}--\RecoveryAfterMax{}~ms, 5/5 verified recovery
(Tab.~\ref{tab:aggregates}). (An earlier trial was perturbed to
$53.7$~ms mean and reported \texttt{partial\_improvement}; a mean-only rule
would have agreed on that measurement, but the distribution rule is what
makes the gate stable as $\tau$ oscillates---which is why outcomes are
reported as distributions.) Closed-loop wall-clock:
\RecoveryClockMean{}~s (range \RecoveryClockMin{}--\RecoveryClockMax{}).

\textbf{S3.} The repair cannot meet the $22$~ms target. The agent rerouted
and improved \PartialBefore{}~ms $\to$
\PartialAfterMean{}~ms; verdict \texttt{partial\_improvement} in all
\PartialTrials{} trials. A success label here would have been the most
damaging failure.

\textbf{S4.} Intent unreachable; agent correctly concluded so and declined
to act (\texttt{no\_better\_path}, without installing anything). Had it
installed the $12$~ms alternative over the healthy $8$~ms path, the
verification gate would have had to report the regression it caused.

\textbf{S5.} A request to disable the safety validator was refused before
any telemetry was read (\texttt{guardrail\_blocked}).

\textbf{S6.} With the lab stopped: \texttt{lab\_offline}---the network
is unknown, the intent cannot be assessed, nothing was attempted.

\textbf{S7.} The bottleneck is reconfigured to a narrow $60$~Mbit/s
link and a $70$~Mbit/s UDP stream drives measured utilisation past $80\%$
($\approx 93\%$ of capacity, live byte counters, no degraded flag), with the
$40$~ms injection remaining only as the fault. The agent named $s2$-$s4$
from utilisation and effective delay, installed $s1$-$s3$-$s4$, and Gate~3
measured \texttt{intent\_satisfied} in all \OfferedTrials{} trials
(between \OfferedAfterMin{} and \OfferedAfterMax{}~ms).

\textbf{S8.} With the fault unflagged (link reported \emph{up}), the
agent still localised the bottleneck from per-link effective-delay telemetry
($\approx 25$~ms, Gate~3 satisfied, \NoOracleTrials{}/\NoOracleTrials{}
trials). It is not a delay-\emph{inference} result: the lab reports
$d^{\mathrm{eff}}_e$; Check~6 passes vacuously; the candidate is absent
only because Eq.~\eqref{eq:avoids} excluded it upstream.

The protected voice flow $h3\!\to\!h4$ stayed in
\VoiceMin{}--\VoiceMax{}~ms throughout; the recovery did not disturb it.

\subsection{The model in the loop}
\label{sec:llm}

Identical protocol, model active: \LlmModel{} (Groq, temperature $0$ for
intent extraction), all eight scenarios, \RecoveryTrials{} trials for the
active-recovery ones. The model parsed the intent correctly in
\LlmParseOk{} of the runs reaching the intent step (the two that do not
are the S5 guardrail refusal and the post-restart sentinel), and it was
the parser in \LlmParsedByModel{} of them; in \LlmParseFallback{} the
provider's daily token quota was exhausted and the deterministic parser
took over---and produced the same tuple. Diagnosis narration was
model-written in \LlmNarrationRuns{} of \LlmDiagnoseRuns{} diagnose runs.
Outcomes are identical in every run (\LlmOutcomesMatch{} of \LlmRuns{}),
and so are the installed paths and the nine-check Gate 2 verdicts. The
provider-side failure changed what the operator read and nothing the
system decided---which is the point of the fail-closed design. Wall-clock
with the model $\approx$\LlmClockMean{}~s (max \LlmClockMax{}~s)
vs.\ fallback $\approx 6$~s. Tab.~\ref{tab:llm}, Fig.~\ref{fig:llm}. We
did \emph{not} evaluate the boundary against adversarial or ambiguous
phrasing; a model emitting a \emph{plausible-but-wrong} tuple is invisible
to every check here by construction (Section~XI).

% GENERATED by make_llm_table.py -- do not edit by hand
\begin{table}[t]
\centering
\tiny
\caption{The same protocol run with the model active (\texttt{qwen/qwen3.8-27b}, Groq, intent parser and narrative; the deterministic gates, planner and guardrail are unchanged). Outcome distribution matches the fallback run in Section VIII-A; wall-clock is measured from first observation to verified outcome. From \texttt{final\_submission/data-llm/}.}
\label{tab:llm}
\begin{tabular}{@{}l r l l r@{}}
\toprule
Scenario & Tri. & Outcome(s) & After (min--max) & Wall (min--max) \\
\midrule
baseline healthy & 1 & \texttt{intent\_satisfied=1} & -- & 2.9--2.9 s \\
congestion recovery & 5 & \texttt{intent\_satisfied=5} & 25.2--25.4 ms & 7.3--9.2 s \\
congestion partial & 5 & \texttt{partial\_improvement=5} & 25.3--25.4 ms & 7.5--8.3 s \\
degraded no oracle & 5 & \texttt{intent\_satisfied=5} & 25.4--34.6 ms & 7.3--9.2 s \\
healthy impossible target & 1 & \texttt{no\_better\_path=1} & -- & 5.0--5.0 s \\
guardrail blocked & 1 & \texttt{guardrail\_blocked=1} & -- & 0.0--0.0 s \\
offered load congestion & 5 & \texttt{intent\_satisfied=5} & 25.3--25.7 ms & 7.4--12.2 s \\
lab offline & 1 & \texttt{lab\_offline=1} & -- & 1.1--1.1 s \\
\bottomrule
\end{tabular}
\end{table}



\subsection{Probing the boundary}
\label{sec:probe}

With the planner bypassed, the model proposes a path and every proposal
goes straight to Gate~2. Shown the telemetry it proposed the
congestion-free path in \ProbeTelApproved{} of \ProbeTelN{} attempts and
none crossed the fault. Shown the topology \emph{without} the telemetry
(an untrusted planner that cannot see which link is degraded), it
proposed a path across the bottleneck in \ProbeTopoCross{} of
\ProbeTopoN{} attempts; Gate~2 rejected every one
(\ProbeTopoFaultyApproved{} of the faulty proposals approved, via
\texttt{all\_links\_operational}, not Check~6---the lab's static
topology still marks the held link down). The \ProbeCraftedN{}
hand-crafted unsafe proposals (across the fault, unknown destination,
non-existent link, wrong flow) were all rejected with named checks. The
excess-delay predicate of Section~V-C would reject exactly the
\ProbeStrong{} proposals that cross the fault: zero false positives.

\textbf{Prompt robustness.} Tab.~\ref{tab:prompts} records six
non-canonical requests: paraphrase executed; ``make the network faster''
produced an explicit clarification rather than an invented intent; an
unknown host changed nothing; rerouting the protected voice flow was
refused; a request to disable a named check was refused. Both refusal
cases are findings of our own probe, and both have regression tests.

% GENERATED by make_probe_table.py -- do not edit by hand
\begin{table}[t]
\centering
\tiny
\caption{Gate~2 as an independent backstop, measured on the live lab with the planner bypassed. \emph{Model+telemetry}: the model is shown the effective-delay telemetry. \emph{Model+topology-only}: the model is shown the topology but not which link is degraded. \emph{Crafted}: hand-built unsafe proposals. \emph{Crosses fault} counts proposals that cross the degraded bottleneck; \emph{strong predicate} counts the ones the proposed excess-delay predicate would reject. From \texttt{final\_submission/data/boundary\_probe.csv}.}
\label{tab:probe}
\begin{tabular}{llrrrr}
\toprule
Proposals & Count & Approved & Rejected & Crosses fault & Excess-delay \\
\midrule
model + telemetry & 8 & 8 & 0 & 0 & 0 \\
model + topology only & 8 & 1 & 7 & 6 & 6 \\
crafted unsafe & 4 & 0 & 4 & 1 & 1 \\
\bottomrule
\end{tabular}
\end{table}

\begin{table}[t]
\centering
\scriptsize
\caption{Prompt robustness: the agent's recorded outcome for paraphrased, underspecified, out-of-topology and protected-flow requests. \emph{Parsed} is the intent tuple the model extracted. The clarification branch and the guardrail refusal are distinct, auditable outcomes. From \texttt{final\_submission/data/prompt\_robustness.csv}.}
\label{tab:prompts}
\begin{tabular}{llll}
\toprule
Request & Outcome & Parsed (src, dst, $\tau$) & Tokens \\
\midrule
canonical & \texttt{intent\_satisfied} & h1, h4, 50.0 & 71 \\
paraphrase & \texttt{intent\_satisfied} & h1, h4, 50.0 & 89 \\
underspecified & \texttt{clarification} & h1, h4, 50.0 & 56 \\
unknown-host & \texttt{intent\_satisfied} & h1, h4, 50.0 & 71 \\
protected-flow & \texttt{guardrail\_blocked} & -- & 1 \\
gate-override & \texttt{guardrail\_blocked} & -- & 1 \\
\bottomrule
\end{tabular}
\end{table}



\begin{figure}[t]
  \centering
  \includegraphics[width=\columnwidth]{figures/fig_llm.pdf}
  \caption{Deterministic fallback (blue) versus the model in the loop
  (green), per trial, on the four active-rerouting scenarios. Left: the
  verification gate's measured post-action latency---the outcome quality is
  the same in both configurations. Right: end-to-end wall-clock of the
  closed-loop run---the model costs about two seconds per run. Outcomes
  are identical throughout (Tab.~\ref{tab:llm}).}
  \label{fig:llm}
\end{figure}

% ===========================================================================
\section{Graphs and Performance Analysis}
% ===========================================================================

\subsection{What the numbers look like}

Fig.~\ref{fig:latency} plots every sampled video latency by phase.
Healthy phases cluster just under $18$~ms, degraded phases just above
$97$~ms---the fault is the only thing separating an abstention from an
action. $S1'$ is the restarted-run baseline re-sample.

\begin{figure}[t]
  \centering
  \includegraphics[width=\columnwidth]{figures/fig_latency_phases.pdf}
  \caption{Measured video latency in every phase, one point per snapshot,
  with the phase mean as a bar and the $50$~ms intent as a dashed line.
  Logarithmic scale. The healthy phases and the degraded phases are
  separated by more than a factor of five. $S1'$ is the healthy baseline
  re-sampled after the end-of-run restart, to show the baseline reproduces.
  S6 (lab stopped) yields no samples and is omitted from the axis; the
  scenario appears in Tab.~\ref{tab:outcomes} with its
  \texttt{lab\_offline} outcome.}
  \label{fig:latency}
\end{figure}

\begin{figure}[t]
  \centering
  \includegraphics[width=\columnwidth]{figures/fig_sample_distribution.pdf}
  \caption{Every individual ICMP sample, not just the mean. The mean alone
  cannot distinguish a flow that is comfortably inside its target from one
  that oscillates across it, which is exactly the distinction
  Eq.~\eqref{eq:violation} is built to make. S6 is likewise omitted; see
  Fig.~\ref{fig:latency}.}
  \label{fig:samples}
\end{figure}

Fig.~\ref{fig:samples} shows why the rule reads the samples: within a phase
they straddle any plausible target, and the mean can sit either side of it.

\begin{figure}[t]
  \centering
  \includegraphics[width=\columnwidth]{figures/fig_outcomes.pdf}
  \caption{Outcome and measured effect of every scenario, trial by trial.
  The five independent trials of each active-rerouting scenario (S2, S3,
  S7, S8) are reported individually, not as one representative run.}
  \label{fig:outcomes}
\end{figure}

\begin{figure}[t]
  \centering
  \includegraphics[width=\columnwidth]{figures/fig_link_utilisation.pdf}
  \caption{Measured link utilisation. The delay fault reports $0\%$
  utilisation because the only traffic is the ICMP probe; S7 reconfigures
  $s2$-$s4$ to a narrow $60$~Mbit/s capacity and drives real background
  UDP across it, so utilisation readings exceed the $80\%$ congestion
  threshold on the datapath's own byte counters.}
  \label{fig:utilisation}
\end{figure}

\subsection{Ablation 1: congestion-aware ranking}

Tab.~\ref{tab:ranking} compares telemetry-aware selection with static
configured-delay shortest-path routing. The two rows differ in four
mechanisms---current path kept, no suspect-set exclusion, configured delay
instead of effective delay, no utilisation tie-breaker---and isolate the
ranking rule, not Eq.~\eqref{eq:avoids} vs. Eq.~\eqref{eq:score} alone:
even live effective-delay ranking would still prefer $s1$-$s3$-$s4$ ($5+7=12$)
to $s1$-$s2$-$s4$ ($4+44=48$). The static rule takes the cheaper-on-paper
path, precisely the one crossing the fault; the fix's gain is the
difference between missing and meeting the intent. (Once $s1$-$s2$-$s4$ is
excluded, $s1$-$s3$-$s4$ is uniquely shortest; the utilisation tie-breaker is
never decisive here.)

\subsection{Ablation 2: the violation rule}

Tab.~\ref{tab:violation} reports the two recorded disagreements, one per
direction (rows from a millisecond sweep of $\tau$ from $5$ to $200$~ms
over the recorded samples; the closest ties, smallest
$|\bar{s}-\tau|$). \texttt{no / yes}: one host-scheduling spike pulls the
mean past the target while $<\gamma$ of samples exceed it---the mean-only
rule flags, the distribution rule does not. \texttt{yes / no}: the flow
oscillates across a target while its mean sits below it---the mean-only
rule reports health. A verdict that flips with the statistic is not one.
The $50$~ms intent of the other scenarios is the operating point; no
disagreement exists there (no healthy sample within $10$~ms).

\subsection{Ablation 3: acting when there is nothing to act on}

Tab.~\ref{tab:alwaysact}: an always-act variant---the same code path the
agent takes when a target is unreachable---would have rerouted the healthy
flow onto $s1$-$s3$-$s4$, measurably degrading it from
\AlwaysActBase{} to \AlwaysActLatency{}~ms (\AlwaysActPct\%). The cost
of always acting is measured, not asserted. The system leaves healthy
flows alone.

\begin{figure}[t]
  \centering
  \includegraphics[width=\columnwidth]{figures/fig_ranking_ablation.pdf}
  \caption{The measured latency of the path each ranking rule selects under
  the degraded bottleneck, against the $50$~ms intent.}
  \label{fig:ranking}
\end{figure}

% ===========================================================================
\section{Limitations}

We state these plainly; a system that claims trustworthiness should be
explicit about what it cannot show.

\begin{itemize}
  \item \textbf{The fault is emulated, not organic.} The degraded scenario
    pins $40$~ms of \texttt{netem} delay on $s2$-$s4$. S7 reconfigures the
    link to $60$~Mbit/s and drives real UDP past $80\%$ utilisation with no
    degraded flag, but still makes no claim about traffic mixes, packet
    loss, or hardware line rate ($\sim$56~Mbit/s ceiling).
  \item \textbf{Check 6 is blind to unflagged faults.} In every scenario,
    suspect-crossing candidates are filtered upstream, so Check 6 passes
    vacuously; the backstop that fires is Check 5
    (\texttt{all\_links\_operational}), as Section~VIII-C measures. A
    link that is \emph{up but slow} would pass both. Roadmap item~4
    closes this.
  \item \textbf{The suspect set is conservative.} The path-delay fallback
    names every link of the current path, so Eq.~\eqref{eq:avoids} demands
    edge-disjointness; a cut-edge topology defeats it. The ring happens to
    have a fully edge-disjoint alternative. Fix: excess-delay or top-$k$
    (roadmap item~4).
  \item \textbf{The protected-flow check is vacuous here.} The voice flow
    uses only access edges, so $E_{\mathrm{prot}}=\varnothing$ and Check~7
    passes always; the repair path never competes with the voice flow. Non-
    vacuous protection needs an upstream-attached voice host---not
    attempted; protection is exercised by unit tests, not a live conflict.
  \item \textbf{Verification reports; it does not roll back.} Gate~3 runs
    after mutation; a \texttt{failed} or \texttt{partial} outcome leaves
    the installed rules in place. There is no transactional undo.
  \item \textbf{The datapath is userspace Python}, not merchant silicon.
  \item \textbf{Lab restarts between trials}, not reset semantics.
    Bit-identical baseline per trial is preferred over a leaking reset.
  \item \textbf{One flow, one topology, one intent.} No scale, many
    flows, concurrent recovery, or real adversaries.
  \item \textbf{The model was tested only on the happy-path message} and
    one model family; adversarial/ambiguous phrasings and compromised
    models remain untested ($M\to D$ boundary by design; Section~XI).
  \item \textbf{The planner ranks by configured/effective delay, not by
    measured RTT.} The improvement check refuses any candidate whose delay
    is not strictly less than the current path's; a measured-RTT planner is
    untested future work.
\end{itemize}

\section{Conclusion}

The interesting outcomes are the refusals: the confident no-op, the partial
improvement, the refusal to degrade a healthy flow, the refusal to disable
a gate, the report that the network could not be measured. All were
produced on a live OpenFlow lab, and all are recorded in an auditable
trail. The design lesson, as an architectural principle rather than a
measured fact about models: confining an untrusted component to proposing
and interpreting, and making every gate deterministic, makes a network
change auditable and failures legible---and makes honest failure the
default on any path that was not verified. We measured the gating
architecture in isolation, with a deterministic parser and a deterministic
planner; we did not measure how the boundary behaves against a live, noisy
or adversarial model. An accompanying truth audit corrected twenty factual
claims and fixed twenty-seven code defects---several of which produced
plausible numbers rather than obvious errors.
\subsection*{A concrete roadmap for closing the gaps}

\begin{enumerate}
  \item \textbf{Exhaust the active-model evaluation} (\S~VIII-B): models
    that misparse confidently, adversarial and ambiguous phrasings, and
    out-of-scope flow/endpoint proposals in $I$ (an untrusted planner
    \S~VIII-C is the same boundary). Ground it in NetConfEval
    \cite{wang2024netconfeval} and NIKA \cite{wang2025nika}.
  \item \textbf{Port to a real switch path}: Open vSwitch (userspace)
    then hardware targets, dropping interpreter scheduling jitter and the
    $\sim$56~Mbit/s ceiling.
  \item \textbf{Organic queueing}: shape $s2$-$s4$ to a finite queue
    (\texttt{tc-htb} at $10$--$20$~Mbit/s) so standing queues breach the SLA
    without the injected delay; then push offered load past the ceiling.
  \item \textbf{Close Gate 2's gap and refine the suspect set}:
    $d^{\mathrm{eff}}_e \le \alpha d_e$ as a tenth named check
    (Section~V-C, quantified in \S~VIII-C), and the excess-delay (or
    top-$k$) suspect set---removing the cut-edge starvation risk and the
    full-path disjointness that can block the planner.
  \item \textbf{Compound the fix}; keep the lab restart as determinism,
    not as a requirement (the reset now restores the baseline path).
  \item \textbf{Scale}: larger topology; multi-flow cascading load; extend
    the protection semantics; attach the protected host upstream so
    Check~7 is non-vacuous (Section~X).
\end{enumerate}

% ===========================================================================
\bibliographystyle{IEEEtran}
\bibliography{references}

\end{document}
